\chapter{검증 결과(4): 통합 회로와 충전 중 보조버스(H3)}

\section{통합 회로 모델: 주행(TB06)}

Stage B·C를 설계할 때의 등가(Stage A: 이상 2:1 전압원 + 경로 저항 13.2 mΩ, Stage C: 정전력 부하)가 실제 회로에서도 성립하는지 확인하려고 세 단을 한 회로로 이었다(배터리 → Stage A 모듈 1 → V\textsubscript{A} 400 µF → Stage B → 400 V 버스 → Stage C 3 kW + 정전력 부하). 세 사례(640 V 15→25 kW, 788 V 5→15 kW, 920 V 25→15 kW)가 모두 통과했고, Stage B 지표는 스위칭 ODE와 일치했으며(버스 차 ≤ 0.38 V, 경로전류 차 ≤ 1.2 A), V\textsubscript{A} 강하는 경로 포함 13.6 mΩ(설계 등가 13.2 mΩ), Stage A 턴오프 전류는 0.06–2.9 \%, 48 V는 ±2 \% 안이었다. TB06·TB07은 150 kHz Stage A로 했고, A25R의 회로 확인은 TB27(13장)과 TB28(12장)이다.


\section{충전 중 보조버스(E16, H3)}

급속충전 중에는 Stage A 4모듈이 모두 동작해 경로저항이 공칭의 42 \%(150 kHz 설계 5.55 mΩ; A25R은 5.25 mΩ)가 되고, 충전기 전류가 V\textsubscript{A}를 움직인다(375 A 변화에 150 kHz 설계 약 0.96 V, A25R 약 0.84 V). 이 조건에서 같은 제어기가 보조버스를 지키는지 평균모델과 스위칭 모델로 시험했고, A25R로 바꾼 뒤 평균모델 시험을 다시 했다(표 \ref{t:e16}, 그림 \ref{f:e16}).

{\footnotesize
\begin{longtable}{L{\dimexpr 0.17\linewidth-2\tabcolsep\relax}L{\dimexpr 0.29\linewidth-2\tabcolsep\relax}L{\dimexpr 0.28\linewidth-2\tabcolsep\relax}L{\dimexpr 0.26\linewidth-2\tabcolsep\relax}}
\caption{충전 중 보조버스 시험 결과(E16, E36)}\label{t:e16}\\
\toprule
\tabhead \textbf{시험} & \textbf{조건} & \textbf{결과(150 kHz Stage A)} & \textbf{A25R 재확인(E36)} \\
\midrule
\endfirsthead
\toprule
\tabhead \textbf{시험} & \textbf{조건} & \textbf{결과(150 kHz Stage A)} & \textbf{A25R 재확인(E36)} \\
\midrule
\endhead
\bottomrule
\endlastfoot
정적 범위 & 640–920 V × 보조부하 −3.2–25 kW, 충전기 명령 I = min[375 A, 150 kW/(V\textsubscript{bat}/2)] & 4,118점 모두 성립. 링크 ≤ 6.28 kW(640 V·25 kW). 632 V 진단점을 더한 4,260점 최대는 링크 6.68 kW·경로 79.9 A & 모두 성립. 충전 중 링크 ≤ 6.29 kW(632 V 포함 6.69 kW) \\
동시 운전 Tier-1 & 460건(632 V 경계 진단 20건, 공차·지연 +3.0 µs 포함). 보조부하 스텝 10가지(0↔10, 5↔15, 10↔20, 15↔25 kW, −3.2↔6.8 kW 회생) & 460/460, 최대 87.43 A(주행 모드 87.62 A) & 460/460(사양 640–920 V만 440/440, E47 재집계) \\
동시 운전 Tier-2 & 78건(632 V 2건 포함, +3 µs 정보 시험 제외). 스텝 1↔25 kW & 78/78, 최대 92.74 A & 78/78(사양 76/76), 최대 92.7 A \\
모드 전환 & 충전기 접속·해제(10 µs, 1 ms), 모듈 기동·셰딩, 240건 & 240/240(Tier-2), Tier-1 전압 품질도 240/240 & 240/240, 경로 ≤ 86.3 A \\
전체 순서 & 주행 → 모듈 기동 → 충전 → +10 kW → 비상 정지 → 셰딩, 30건 & 30/30, 최대 87.04 A, 버스 397.8–401.6 V & 30/30, 버스 397.81–401.44 V \\
스위칭 모델 & 충전 중 Tier-1·Tier-2 40건, 접속·비상 정지 8건 & 48/48, 순시 최대 96.26 A & 다시 하지 않음(회로 TB27·TB28로 확인) \\
\end{longtable}
}
\srcnote{주: A25R 재확인은 Stage A 모듈 출력저항 R\textsubscript{A}를 10.2 mΩ에서 9.0 mΩ로 바꾸고 같은 기준으로 계산했다. 이 평균모델은 V\textsubscript{A}를 V\textsubscript{bat}/2 + R\textsubscript{A}·(i\textsubscript{ch} − i\textsubscript{p})/모듈 수의 정적 관계로 두며, C\textsubscript{A} 2.4 mF와 Stage A 공진·게이트의 결합이 없어 이 근사가 모든 판정량에서 보수적이라는 증명은 없다. A25R의 결합 동특성은 대표 사례(E35, TB27·TB28)에서만 확인했다(18장). 자료: 검증 근거 E16, E36.}

\begin{figure}[!htb]
\centering
\includegraphics[width=15.5cm]{\repo/validation/2026-10-22-charging-integration/E16/E16_charging_transients.png}
\caption{충전 중 부하 스텝과 충전기 접속·해제 응답(평균모델)}\label{f:e16}
\srcnote{주: 150 kHz 설계의 Stage A 등가(R\textsubscript{A} 10.2 mΩ)로 계산한 응답이다. A25R 등가로 다시 계산한 판정은 표 \ref{t:e16}에 있다.}
\end{figure}

\section{충전 통합 회로 모델(TB07)}

TB06에 Stage A 모듈 2–4(8위상)와 400 V급 충전기(제어 전류원)를 더해 기본 3사례(640 V 충전 중 15→25 kW, 920 V 충전 비상 정지, 788 V 충전기 접속)와 추가 4실행(변화 시간 100 µs·1 ms, 640 V·25 kW 접속·정지)을 계산했다. \textbf{회로 한계는 7/7 통과}했고, Stage A 8위상은 역방향 대전류에서도 턴오프 전류 ≤ 1.1 A, 전류 분담 ≤ 0.9 \%를 유지했다.

검증 모델과의 일치는 조건부다. 100 µs 이상의 충전전류 변화는 일치했지만, 10 µs 변화 3개 실행은 경로전류 첨두가 3.3–13.4 A 커서 일치 기준을 넘었다. 원인은 V\textsubscript{A} 노드 공진(주기 약 30 µs)이다. 가장 불리한 640 V·25 kW 접속에서 회로 경로전류는 98.49 A(한계 100 A)였다. 그래서 150 kHz 설계에서는 \textbf{충전전류 증가 램프 ≥ 100 µs}(375 A 기준 ≤ 3.75 A/µs)를 설계 요구로 두었다.

A25R에서는 공진주파수가 25 kHz로 낮아져 100 µs 램프에서 영전류 턴오프 위반이 12–22번 생겼고, 사전 기동 없이 충전을 시작하면 위상 첨두가 217–221 A로 정격을 크게 넘었다(13장). 그래서 최종 요구는 \textbf{충전기 전류가 오르기 전에 Stage A 모듈을 모두 켜는 사전 기동}과 \textbf{충전전류 증가 램프 ≥ 600 µs}(375 A 기준 ≤ 0.625 A/µs)다. 600 µs는 25 kHz 공진 15주기로, 150 kHz의 100 µs와 같은 주기 수다. A25R 회로에서는 느린 충전 시작과 모듈 고장(TB27, 4/4), 충전 중 정지·느린 충전 시작과 종료(TB28, 4/4; 충전 시작·종료 D28은 보조 15 kW, 경로 상실 A28·B28과 모듈 고장 C28은 보조 25 kW·이전 단일 1 Ω 클램프)로 이 요구 아래의 동작을 확인했다.

\textbf{H3 판정: 충족(조건: 충전 전 사전 기동과 충전전류 증가 램프 ≥ 600 µs, 보조부하 25 kW에서 판정).} 600 µs는 명시한 부하(15·25 kW)와 제어 상태에서 확인한 기동 조건이다. 같은 기동·램프에서 640 V·보조부하 1 kW는 위상 첨두 78.74 A로 E35 공칭 한도 74.53 A를 넘었고(0 kW 78.63 A, 적분 간격 절반에서도 78.68 A, 턴오프 위반 0, E46), 경부하를 포함한 전 범위의 충분조건으로 쓰려면 버스 부하·전압·초기 모듈 수·공차·램프를 교차한 기동/감발 규칙을 따로 검증해야 한다. \textbf{H1 충전 부분: 충족.} 회로에서도 충전전류는 V\textsubscript{A} 노드에서 Stage A로 바로 들어가고 경로전류는 보조부하만큼만 흐른다(375 A 충전 중 13–79 A). 충전 중 링크전력은 사양 안에서 정적 6.29 kW·과도 정식 7.04 kW로 8 kW 안이다(632 V 경계 진단 6.69·7.49 kW; 7.49 kW 사례는 경로전류가 93.75 A를 넘은 정보용 사례). 램프 요구는 차량 쪽 요청 램프로 만족시킬 수 있지만 충전기마다 응답 속도가 달라 하드웨어 단계에서 확인해야 한다(19장).
